% (c) Nikita Lisitsa, lisyarus@gmail.com, 2026

\documentclass[10pt]{beamer}

\usepackage[T2A]{fontenc}
\usepackage[russian]{babel}
\usepackage{minted}

\usepackage[most]{tcolorbox}
\tcbuselibrary{minted}

\usepackage{graphicx}
\graphicspath{ {./images/} }

\usepackage{adjustbox}

\usepackage{color}
\usepackage{soul}
\usepackage{multicol}
\usepackage{multirow}
\usepackage{hyperref}
\usepackage{tikz}
\usetikzlibrary{arrows.meta,positioning,calc}

\usetheme{metropolis}

\definecolor{red}{rgb}{1,0,0}
\definecolor{green}{rgb}{0,0.5,0}
\definecolor{codebg}{RGB}{29,35,49}
\setminted{fontsize=\footnotesize}

\makeatletter
\newcommand{\slideimage}[1]{
  \begin{figure}
    \begin{adjustbox}{width=\textwidth, totalheight=\textheight-2\baselineskip-2\baselineskip,keepaspectratio}
      \includegraphics{#1}
    \end{adjustbox}
  \end{figure}
}
\makeatother

\title{Язык программирования C}
\subtitle{Лекция 8: Макросы}
\author{Лисица Никита}
\date{2026}

\setbeamertemplate{footline}[frame number]

\usemintedstyle{tango}

\begin{document}

\frame{\titlepage}

\begin{frame}[fragile]
\frametitle{Макросы}
\begin{itemize}
\item Макросы -- инструмент текстовой замены одних выражений на другие
\pause
\item Выполняются \textbf{препроцессором}, до основного прохода компилятора -- т.е. до проверки типов, до разрешения имён, и т.п.
\pause
\item Заменить можно \textbf{не любое} выражение: макросы всегда выглядят, как переменные или как функции:
\end{itemize}
\usemintedstyle{lightbulb}
\begin{tcblisting}{listing engine=minted, minted language=cpp, minted options={fontsize=\scriptsize}, listing only, colback=codebg, colframe=codebg, rounded corners}
#define PI 3.1415926535
#define SUM(a, b) a + b
\end{tcblisting}
\pause
\begin{multicols}{2}
\begin{tcblisting}{listing engine=minted, minted language=cpp, minted options={fontsize=\scriptsize}, listing only, colback=codebg, colframe=codebg, rounded corners}
// До препроцессора:

double x = sin(PI);
int s = SUM(a, b);
\end{tcblisting}
\columnbreak
\begin{tcblisting}{listing engine=minted, minted language=cpp, minted options={fontsize=\scriptsize}, listing only, colback=codebg, colframe=codebg, rounded corners}
// После препроцессора:

double x = sin(3.1415926535);
int s = a + b;
\end{tcblisting}
\end{multicols}
\end{frame}

\begin{frame}[fragile]
\frametitle{Макросы}
\begin{itemize}
\item Если результат раскрытия макроса содержит макросы, они тоже раскрываются (но рекурсия запрещена):
\end{itemize}
\usemintedstyle{lightbulb}
\begin{tcblisting}{listing engine=minted, minted language=cpp, minted options={fontsize=\scriptsize}, listing only, colback=codebg, colframe=codebg, rounded corners}
#define PI 3.1415926535
#define SIN_DEG(x) sin(x / 180.0 * PI)
#define MY_ANGLE 30.0

double x = SIN_DEG(MY_ANGLE);

// Превратится в
double x = 30.0 / 180.0 * 3.1415926535;
\end{tcblisting}
\end{frame}

\begin{frame}[fragile]
\frametitle{Макросы: if}
\begin{itemize}
\item Макросы можно использовать для условного включения кода:
\end{itemize}
\usemintedstyle{lightbulb}
\begin{tcblisting}{listing engine=minted, minted language=cpp, minted options={fontsize=\scriptsize}, listing only, colback=codebg, colframe=codebg, rounded corners}
#define USE_SIMD

// ...

#ifdef USE_SIND
  #include <immintrin.h>
#endif

vec4 add(vec4 a, vec4 b) {
  #ifdef USE_SIMD
    __m128 va = _mm_loadu_ps(&a.x);
    __m128 vb = _mm_loadu_ps(&b.x);
    __m128 vr = _mm_add_ps(va, vb);
    vec4 result;
    _mm_storeu_ps(&result.x, vr);
    return result;
  #else
    return (vec4){a.x + b.x, a.y + b.y, a.z + b.z, a.w + b.w};
  #endif
}
\end{tcblisting}
\end{frame}

\begin{frame}[fragile]
\frametitle{Макросы: if}
\begin{itemize}
\item \mintinline{cpp}|#if| может содержать более сложные арифметические и логические выражения:
\end{itemize}
\usemintedstyle{lightbulb}
\begin{tcblisting}{listing engine=minted, minted language=cpp, minted options={fontsize=\scriptsize}, listing only, colback=codebg, colframe=codebg, rounded corners}
#if defined(__APPLE__) && LIB_VERSION >= 3
  ...
#endif
\end{tcblisting}
\pause
\begin{itemize}
\item Полезно для версионирования кода, определения окружения компилятора (ОС, битность системы), и т.п.
\end{itemize}
\end{frame}

\begin{frame}[fragile]
\frametitle{Макросы: проблемы}
\begin{itemize}
\item Так как это просто текстовая подстановка (игнорирующая синтаксис C), в их использовании легко сделать ошибку
\end{itemize}
\pause
\usemintedstyle{lightbulb}
\begin{tcblisting}{listing engine=minted, minted language=cpp, minted options={fontsize=\scriptsize}, listing only, colback=codebg, colframe=codebg, rounded corners}
#define MUL(a, b) a * b

int x = 5;
int y = MUL(x + 2, 7);

printf("%i", y);
\end{tcblisting}
\pause
\usemintedstyle{tango}
\begin{itemize}
\item Что выведется на экран? \pause 19, так как код превратится в \mintinline{cpp}|x + 2 * 7|
\end{itemize}
\end{frame}

\begin{frame}[fragile]
\frametitle{Макросы: проблемы}
\begin{itemize}
\item Как это исправить?
\end{itemize}
\pause
\usemintedstyle{lightbulb}
\begin{tcblisting}{listing engine=minted, minted language=cpp, minted options={fontsize=\scriptsize}, listing only, colback=codebg, colframe=codebg, rounded corners}
#define MUL(a, b) (a) * (b)

int x = 5;
int y = MUL(x + 2, 7);

printf("%i", y);
\end{tcblisting}
\pause
\usemintedstyle{tango}
\begin{itemize}
\item Теперь выведется 49, так как код превратится в \mintinline{cpp}|(x + 2) * (7)|
\end{itemize}
\end{frame}

\begin{frame}[fragile]
\frametitle{Макросы: проблемы}
\usemintedstyle{lightbulb}
\begin{tcblisting}{listing engine=minted, minted language=cpp, minted options={fontsize=\scriptsize}, listing only, colback=codebg, colframe=codebg, rounded corners}
#define ABS(x) ((x) >= 0 ? (x) : -(x))

int x = ABS(very_slow_function());

printf("%i", x);
\end{tcblisting}
\pause
\usemintedstyle{tango}
\begin{itemize}
\item В чём проблема? \pause \mintinline{cpp}|very_slow_function()| вызовется два раза
\end{itemize}
\end{frame}

\begin{frame}[fragile]
\frametitle{Макросы: проблемы}
\usemintedstyle{lightbulb}
\begin{tcblisting}{listing engine=minted, minted language=cpp, minted options={fontsize=\scriptsize}, listing only, colback=codebg, colframe=codebg, rounded corners}
#define ABS(x) ((x) >= 0 ? (x) : -(x))

int x = 10;
int y = ABS(x++);

printf("%i %i", x, y);
\end{tcblisting}
\pause
\usemintedstyle{tango}
\begin{itemize}
\item Что выведется на экран? \pause \mintinline{text}|12 11| (инкремент применится дважды)
\pause
\item \textbf{NB:} До C11 это UB, так как выражение содержит несколько побочных эффектов для одной и той же переменной
\end{itemize}
\end{frame}

\begin{frame}[fragile]
\frametitle{Макросы: проблемы}
\begin{itemize}
\item Как исправить \mintinline{cpp}|ABS(x)|?
\pause
\item Никак: лучше написать функцию
\end{itemize}
\usemintedstyle{lightbulb}
\begin{tcblisting}{listing engine=minted, minted language=cpp, minted options={fontsize=\scriptsize}, listing only, colback=codebg, colframe=codebg, rounded corners}
inline int abs(int x) {
  return x >= 0 ? x : -x;
}
\end{tcblisting}
\end{frame}

\begin{frame}[fragile]
\frametitle{Макросы: проблемы}
\usemintedstyle{lightbulb}
\begin{tcblisting}{listing engine=minted, minted language=cpp, minted options={fontsize=\scriptsize}, listing only, colback=codebg, colframe=codebg, rounded corners}
#define SWAP(a, b) int temp = a; a = b; b = temp;
\end{tcblisting}
\pause
\usemintedstyle{tango}
\begin{itemize}
\item В чём проблема? \pause Нельзя использовать в функции, где уже объявлена переменная \mintinline{cpp}|temp| (будет multiple definition):
\end{itemize}
\usemintedstyle{lightbulb}
\begin{tcblisting}{listing engine=minted, minted language=cpp, minted options={fontsize=\scriptsize}, listing only, colback=codebg, colframe=codebg, rounded corners}
int foo(int x) {
  int temp = x * x;
  SWAP(x, temp);
  return x;
}
\end{tcblisting}
\end{frame}

\begin{frame}[fragile]
\frametitle{Макросы: проблемы}
\begin{itemize}
\item Как это исправить? \pause Явно создать вложенный скоуп (но лучше всё-таки функция):
\end{itemize}
\usemintedstyle{lightbulb}
\begin{tcblisting}{listing engine=minted, minted language=cpp, minted options={fontsize=\scriptsize}, listing only, colback=codebg, colframe=codebg, rounded corners}
#define SWAP(a, b) { int temp = a; a = b; b = temp; }
\end{tcblisting}
\pause
\begin{itemize}
\item Начиная с C23 можно не указывать тип:
\end{itemize}
\usemintedstyle{lightbulb}
\begin{tcblisting}{listing engine=minted, minted language=cpp, minted options={fontsize=\scriptsize}, listing only, colback=codebg, colframe=codebg, rounded corners}
#define SWAP(a, b) { typeof(a) temp = a; a = b; b = temp; }
\end{tcblisting}
\end{frame}

\begin{frame}[fragile]
\frametitle{Многострочные макросы}
\begin{itemize}
\item Можно писать макросы в несколько строк, для этого в конце каждой строки кроме последней нужно поставить обратный слеш:
\end{itemize}
\usemintedstyle{lightbulb}
\begin{tcblisting}{listing engine=minted, minted language=cpp, minted options={fontsize=\scriptsize}, listing only, colback=codebg, colframe=codebg, rounded corners}
#define SWAP(a, b) { \
  typeof(a) temp = a; \
  a = b; \
  b = temp; }
\end{tcblisting}
\end{frame}

\begin{frame}[fragile]
\frametitle{Макросы: stringize}
\begin{itemize}
\item Макросы позволяют превратить любое выражение в строку с помощью оператора препроцессора \mintinline{cpp}|#|:
\end{itemize}
\usemintedstyle{lightbulb}
\begin{tcblisting}{listing engine=minted, minted language=cpp, minted options={fontsize=\scriptsize}, listing only, colback=codebg, colframe=codebg, rounded corners}
#define STR(x) #x

printf(STR(x + 10));

// Превратится в
printf("x + 10");
\end{tcblisting}
\pause
\usemintedstyle{tango}
\begin{itemize}
\item Раскрытия макросов внутри \mintinline{cpp}|x| при этом \textbf{не происходит}!
\end{itemize}
\end{frame}

\begin{frame}[fragile]
\frametitle{Макросы: stringize}
\begin{itemize}
\item Полезно для логирования, assert'ов, unit-тестов, и т.п.:
\end{itemize}
\usemintedstyle{lightbulb}
\begin{tcblisting}{listing engine=minted, minted language=cpp, minted options={fontsize=\scriptsize}, listing only, colback=codebg, colframe=codebg, rounded corners}
#define ASSERT(x) if(!(x)) { \
    printf("Assertion failed: " #x "\n"); \
    exit(1); \
  }

int x = 10;
ASSERT(x < 0);
\end{tcblisting}
\begin{itemize}
\item Выведет:
\end{itemize}
\usemintedstyle{lightbulb}
\begin{tcblisting}{listing engine=minted, minted language=text, minted options={fontsize=\scriptsize, formatcom=\color{white}}, listing only, colback=codebg, colframe=codebg, rounded corners}
Assertion failed: x < 0
\end{tcblisting}
\end{frame}

\begin{frame}[fragile]
\frametitle{Макросы: конкатенация}
\begin{itemize}
\item Можно склеивать несколько выражений в одно с помощью оператора препроцессора \mintinline{cpp}|##|:
\end{itemize}
\usemintedstyle{lightbulb}
\begin{tcblisting}{listing engine=minted, minted language=cpp, minted options={fontsize=\scriptsize}, listing only, colback=codebg, colframe=codebg, rounded corners}
#define CONCAT(x, y) x##y

CONCAT(pri, ntf)("Hello, world!");

// Превратится в
printf("Hello, world!");
\end{tcblisting}
\pause
\usemintedstyle{tango}
\begin{itemize}
\item Раскрытия макросов внутри \mintinline{cpp}|x| и \mintinline{cpp}|y| при этом тоже \textbf{не происходит}!
\end{itemize}
\end{frame}

\begin{frame}[fragile]
\frametitle{Макросы: кодогенерация}
\begin{itemize}
\item Конкатенацию можно использовать для генерации \textit{generic} кода (для произвольных типов данных):
\end{itemize}
\usemintedstyle{lightbulb}
\begin{tcblisting}{listing engine=minted, minted language=cpp, minted options={fontsize=\scriptsize}, listing only, colback=codebg, colframe=codebg, rounded corners}
#define MAKE_MAX(type) type max_##type(type a, type b) { \
  return a > b ? a : b; }

MAKE_MAX(int)    // Сгенерирует max_int
MAKE_MAX(float)  // Сгенерирует max_float
MAKE_MAX(double) // Сгенерирует max_double

int x = max_int(10, 11);
double y = max_double(3.14, 2.71);
\end{tcblisting}
\end{frame}

\begin{frame}[fragile]
\frametitle{Макросы: кодогенерация}
\usemintedstyle{lightbulb}
\begin{tcblisting}{listing engine=minted, minted language=cpp, minted options={fontsize=\scriptsize}, listing only, colback=codebg, colframe=codebg, rounded corners}
#define MAKE_ARRAY(type) typedef struct { \
  type * data; \
  size_t size; \
  } array_##type;

MAKE_ARRAY(int)

array_int a;
a.size = 10;
a.data = malloc(a.size * sizeof(*a.data));
\end{tcblisting}
\end{frame}

\begin{frame}[fragile]
\frametitle{Макросы: кодогенерация}
\usemintedstyle{lightbulb}
\begin{tcblisting}{listing engine=minted, minted language=cpp, minted options={fontsize=\scriptsize}, listing only, colback=codebg, colframe=codebg, rounded corners}
#define MAKE_LIST(type) typedef struct list_##type##_t { \
  type value; \
  struct list_##type##_t * next; \
  } list_##type;

MAKE_LIST(int)

list_int l;
l.value = 10;
l.next = nullptr;
\end{tcblisting}
\end{frame}

\begin{frame}[fragile]
\frametitle{Макросы: упрощение кода}
\usemintedstyle{lightbulb}
\begin{tcblisting}{listing engine=minted, minted language=cpp, minted options={fontsize=\scriptsize}, listing only, colback=codebg, colframe=codebg, rounded corners}
#define FOREACH(size, var) \
  for (int var = 0; var < size; ++var)

FOREACH(10, i) {
  printf("%d\n", i * i);
}
\end{tcblisting}
\pause
\usemintedstyle{lightbulb}
\begin{tcblisting}{listing engine=minted, minted language=cpp, minted options={fontsize=\scriptsize}, listing only, colback=codebg, colframe=codebg, rounded corners}
#define LIST_FOREACH(list, var) \
  for (node * var = list; var != NULL; var = var->next)

LIST_FOREACH(mylist, n)
  printf("%d\n", n->value);
\end{tcblisting}
\end{frame}

\begin{frame}[fragile]
\frametitle{Макросы: логирование}
\begin{itemize}
\item Есть особые нестандартные, но широко распространённые макросы для получения информации о текущем месте в коде
\pause
\item \mintinline{cpp}|__FILE__| раскрывается в имя файла, содержащего этот макрос
\pause
\item \mintinline{cpp}|__LINE__| раскрывается в номер строки в файле
\end{itemize}
\pause
\usemintedstyle{lightbulb}
\begin{tcblisting}{listing engine=minted, minted language=cpp, minted options={fontsize=\scriptsize}, listing only, colback=codebg, colframe=codebg, rounded corners}
#define LOG(msg) printf("(%s:%d) " msg, __FILE__, __LINE__)

// В main.cpp, на 25-ой строке:
LOG("test");
\end{tcblisting}
\begin{itemize}
\item Выведет
\end{itemize}
\usemintedstyle{lightbulb}
\begin{tcblisting}{listing engine=minted, minted language=text, minted options={fontsize=\scriptsize, formatcom=\color{white}}, listing only, colback=codebg, colframe=codebg, rounded corners}
(main.cpp:25) test
\end{tcblisting}
\end{frame}

\begin{frame}[fragile]
\frametitle{Макросы: variadic}
\begin{itemize}
\item Макросы тоже могут быть variadic, т.е. иметь произвольное количество аргументов (обозначается тоже как многоточие \mintinline{cpp}|...|)
\pause
\item \mintinline{cpp}|__VA_ARGS__| -- макрос, раскрывающийся в список всех variadic аргументов (вместе с запятыми между ними)
\end{itemize}
\pause
\usemintedstyle{lightbulb}
\begin{tcblisting}{listing engine=minted, minted language=cpp, minted options={fontsize=\scriptsize}, listing only, colback=codebg, colframe=codebg, rounded corners}
#define LOG(format, ...) printf(format, __VA_ARGS__)

LOG("%s %d", "Hello ", 10);
\end{tcblisting}
\begin{itemize}
\item Выведет
\end{itemize}
\usemintedstyle{lightbulb}
\begin{tcblisting}{listing engine=minted, minted language=text, minted options={fontsize=\scriptsize, formatcom=\color{white}}, listing only, colback=codebg, colframe=codebg, rounded corners}
Hello 10
\end{tcblisting}
\end{frame}

\begin{frame}[fragile]
\frametitle{Макросы: variadic}
\begin{itemize}
\item Если в variadic макрос не передать ни одного variadic аргумента, \mintinline{cpp}|__VA_ARGS__| раскрывается в пустоту
\pause
\item Часто это становится проблемой, если передавать \mintinline{cpp}|__VA_ARGS__| аргументами в какую-то функцию -- в конце остаётся лишняя запятая:
\end{itemize}
\pause
\usemintedstyle{lightbulb}
\begin{tcblisting}{listing engine=minted, minted language=cpp, minted options={fontsize=\scriptsize}, listing only, colback=codebg, colframe=codebg, rounded corners}
#define LOG(format, ...) printf(format, __VA_ARGS__)

LOG("Hello");

// Превратится в
printf("Hello", );
\end{tcblisting}
\end{frame}

\begin{frame}[fragile]
\frametitle{Макросы: variadic}
\begin{itemize}
\item Решение: \mintinline{cpp}|__VA_OPT__| (C23, C++20)
\pause
\item \mintinline{cpp}|__VA_OPT__(X)| раскрывается в \mintinline{cpp}|X|, если список variadic аргументов непустой, иначе раскрывается в пустоту
\end{itemize}
\pause
\usemintedstyle{lightbulb}
\begin{tcblisting}{listing engine=minted, minted language=cpp, minted options={fontsize=\scriptsize}, listing only, colback=codebg, colframe=codebg, rounded corners}
#define LOG(format, ...) \
  printf(format __VA_OPT_(,) __VA_ARGS__)
\end{tcblisting}
\end{frame}

\begin{frame}[fragile]
\frametitle{Макросы: логирование с форматированием}
\usemintedstyle{lightbulb}
\begin{tcblisting}{listing engine=minted, minted language=cpp, minted options={fontsize=\scriptsize}, listing only, colback=codebg, colframe=codebg, rounded corners}
#define LOG(format, ...) \
  printf("(%s:%d) " format, __FILE__, __LINE__ \
    __VA_OPT__(,) __VA_ARGS__)

int a = 10, b = 20;
// В main.cpp, в 25-ой строке
LOG("%d + %d = %d", a, b);
\end{tcblisting}
\begin{itemize}
\item Выведет
\end{itemize}
\usemintedstyle{lightbulb}
\begin{tcblisting}{listing engine=minted, minted language=text, minted options={fontsize=\scriptsize, formatcom=\color{white}}, listing only, colback=codebg, colframe=codebg, rounded corners}
(main.cpp:25) 10 + 20 = 30
\end{tcblisting}
\end{frame}

\end{document}
